% a02_Kpos part deg2: the Ward rows in the R > p branch (the cut by clusters), and the conclusion for (DEG) (input after a02_deg1)
\begin{theorem}[the Ward rows, $R\le p$]\label{thm:a02:WR}
Let the localisation be by centres, and let $S$ be a Ward-row prefix component at a fixed tuple in the branch $R\le p$ of \cite[(3.39)]{K1}, after King's reduction \cite[(3.43)--(3.55)]{K1} with the gauge $B=A_j(x_0)$ of the cube. Then the sum over $F_S$ and $A_S$ is a sum of terms of three shapes: \textup{(R0)} zero; \textup{(R1)} local terms $\sum_x5^{-3j}c_a(x)\prod\partial^a\tilde\Phi(x)$, $|a|\le n=-D(S)$, with $|c_a(x)|\le C\rho_{j_0}^{D(S)+|a|}$ uniformly in $x$, $j_0$, $k$, $\eta$ and the tuple, times $W^\circ$ and the frozen $\tilde A(x)$-legs; \textup{(R2)} remainders with a credit $\rho^\theta$ on a line of $S$ and a debit of order $\le\theta$ on one leg of $S$ or on a unit-scale weight, $D'(S)>0$ and $D'(\gamma)\ge D(\gamma)$ for every prefix component $\gamma\subsetneq S$. The constant depends on $\nbar$, $d$, $L$, $\alpha$ only.
\end{theorem}
\begin{proof}
(1) By \cref{lem:a02:Lsharp}(c) every member of $F_S$, and every member of every family met in the nesting, is reduced with the same $\square$, $\square'$, $x_0$, $B$, $\tilde A$ and lies in the same case below. (2) Insert at every $\tilde A$-vertex $v$ the identity $1=\zeta(v)+(1-\zeta(v))$, $\zeta(v)\defeq\psi(\dist(v,\square)/p)$, $\psi\in C^\infty$, $\psi=1$ on $[0,1]$, $0$ on $[\frac32,\infty)$, $|\psi'|\le4$. An $\tilde A$-vertex, an insertion vertex of the expansion \cite[(3.52)--(3.53)]{K1} of one propagator in $\tilde A$, has exactly two scalar line-ends, both on the chain of that propagator, and external $\tilde A$-legs only; so a connected subgraph all of whose vertices are $\tilde A$-vertices is a path of $r$ vertices inside one chain, with $E_S+E_{S_D}=2$, $E_A\ge r$ and order $\ge r$, hence $D\ge\frac r2-3+\frac12(2-E_{S_D}+r)+\frac32E_{S_D}=r-2+E_{S_D}$ by the degree formula of \cref{ssec:a02:ward}, and $D\le0$ only for $r=2$, $E_{S_D}=0$, $E_A=2$: the tree \cite[(2.4)]{Ba3}, which is not a Ward row (its member-wise identity at every tuple, the cut of its position sums at $\partial\square'$ included, is \cref{prop:a02:TREE} below). So $S$ contains a vertex $u$ of the original graph, localised by its centre within $\sqrt d$ of $\square$. A term in which some $\tilde A$-vertex $v$ carries $1-\zeta$ has $\dist(v,\square)\ge p$, hence $|u-v|\ge p-\sqrt d$; a term with $\dist(x,\square)>p/2$ has $|x-u|\ge p/2-\sqrt d$; a term with a line $\delta G_{(i)}(\square')$ has either both ends of that line at distance $\ge p/2$ from $\partial\square'$, or an end $w$ with $\dist(w,\square)\ge\frac32p$, hence $|w-u|\ge\frac32p-\sqrt d$. Two vertices of $S$ at distance $\ge t$ are joined by a path of at most $V-1$ lines of $S$, $V\le C(\nbar)$ its number of vertices, so some line $\ell$ of $S$, of index $i$, has length $\ge t/V$, and $e^{-\delta_0|\ell|/\rho_i}\le C_\theta(V\rho_i/(\delta_0t))^\theta e^{-\delta_0|\ell|/(2\rho_i)}$ for every $\theta\ge0$ (the configuration space is split by the first line of the path of length $\ge t/V$; the split is sharp but harmless, since no identity is used on these terms); a $\delta G_{(i)}(\square')$ line with both ends at distance $\ge p/2$ from $\partial\square'$ carries $e^{-\delta_1p/(2\rho_i)}$ by \cref{lem:a02:I}. In every case the line supplies a credit of any order, $\le C_\theta(\rho_i/p)^\theta$ with $\theta=3>-D(S)$, on a line of $S$ with no debit and no identity (\cref{prop:a02:FL}, \cref{lem:a02:I}); these terms are of shape (R2), bounded with absolute values, and $S$ needs no rewriting. (3) In the remaining \emph{bulk} case ($\dist(x,\square)\le p/2$, every line a free slice, every $\tilde A$-vertex carrying $\zeta$, so $\zeta(x)=1$ and all position sums are sums over the whole lattice) \cref{lem:a02:FW} turns the weighted family into $W^\circ$ times the unweighted family against $\tilde\Phi$ plus (R2) weight units, and the unweighted family is the zero-background free-lattice object of \cref{lem:a02:RS}: resum and replace by $C^\xi$; for the photon two-point classes $F^C_S$ is the full fixed-order connected two-photon function of the $C^\xi$ theory, all ways of filling the inside and the $\dm$ units included \cite[p.~431]{Ba3}, so $c_0(F^C_S)=0$ by \cref{lem:a02:GC}(a), and at order $e^2$ also $c_1(F^C_S)=0$ by \cref{lem:a02:REF}(b); for the insertion classes the inside families are exactly the insertion families of the order-$2$ parents $\Sigma_1$, the $\lambda$-tadpole and the $Q_2$-tadpole, so $c_0(F^C_S)=-eq\times(\text{first moment of the parent})^C=0$ by \cref{lem:a02:GC}(b), \cref{lem:a02:REF}(a) and the locality of the tadpoles; \cref{lem:a02:RS} bounds the remaining terms by $C\rho_{j_0}^{D+|a|}$. The transport and Taylor steps \cite[(3.10), (3.34)]{Ba3}, with the reference vertex the vertex of the coarsest leg and the Taylor leg the finer one, then give (R1) and (R2) with the degree arithmetic of \cite[Sect.~3]{Ba3}: each remainder factor $|x_e-x|^\theta$ is bounded along a path of lines of $S$ at the price $\rho_{j_l}^\theta$, the H\"older quotient of the leg costs $\rho_{j_e}^{-\theta}$ \cite[(3.65), (3.85)]{K1}, $D'(S)=D(S)+\theta>0$, and a proper prefix component does not contain the debited leg. (4) Sub-units are rewritten first, finest first, at the same tuple; each rewriting is an identity on its own family, so the moments of $S$ are those of its original weighted family kernel, to which (3) applies; a sub-unit's frozen weights are frozen again at $x$, since products of $[0,1]$-valued unit-scale factors compose; each leg receives at most one unit's debit, of order $\le1+\alpha$; a weight debit is never a leg debit; and one gauge, $B=A_j(x_0)$ per cube, is used throughout.
\end{proof}
% a02_Kpos part deg2, continued: the tree near the face of the reduction cube (input from a02_deg2.tex after thm:a02:WR)
\paragraph{The tree near the face of the reduction cube.} Step (2) of the proof of \cref{thm:a02:WR} shows that the only source-free connected subgraph with $D\le0$ all of whose vertices are $\tilde A$-vertices is the tree $T$ of \cite[(2.4)]{Ba3}: two insertion vertices of \cite[(3.52)--(3.53)]{K1} of $J_1$ type joined by a doubly differentiated line, $D(T)=0$, no partner, inside family $\{T\}$. Its member-wise identity, a derivative at the end of a line (\cref{ssec:a02:ward}), is a sum over the whole lattice; at a fixed tuple the positions of insertion vertices are summed over the reduction cube $\square'$ with a sharp cut (\cite[p.~668]{K1}: only through the expansion \cite[(3.52)]{K1} can vertices approach $\partial\square'$), so near $\partial\square'$ the identity leaves a face term. The lines are decomposed scale by scale on $\square'$, $G_{(i)}(\square',0)=G_{(i)}(0)+\delta G_{(i)}(\square')$, the difference kept as an internal line of index $i$ (as in step (2) of the proof of \cref{thm:a02:WR}); we write $F_i(v)\defeq\exp[-\delta_1\dist(v,\partial\square')/\rho_i]$ with $\delta_1\defeq\delta_0/2$, $\delta_0$ the decay rate of the free slices \cite[(3.63)]{K1}.

\begin{lemma}[face factors]\label{lem:a02:FACE}
\textup{(a)} For $u,v\in\square'$ the line $\delta G_{(i)}(\square')(u,v)$ obeys the free slice bounds \cite[(3.63)--(3.65)]{K1} (sup, first differences, H\"older quotients of the first differences) at index $i$ with half the rate, times $F_i(u)F_i(v)$ (for a H\"older pair $u,\tilde u$, $F_i$ taken at the point of the pair nearer to $\partial\square'$). \textup{(b)} Let a term of the ordered bound of \cref{thm:a02:S3}(a) carry the factor $F_i(v)$ at a vertex $v$ of the tree $T$, $i$ the index of its line $\ell$. Then the term is an \textup{(R2)} unit in the sense of \cref{thm:a02:WR}: a credit $\rho_i$ on $\ell$ and a debit $\rho_{i'}^{-1}$ on one leg $l'$ of $T$, of index $i'\ge i$, or no debit when the surviving vertex of $T$ is the root of the graph; $D'(T)=1$, and $D'(\gamma)\ge D(\gamma)$ for every prefix component $\gamma\supsetneq T$.
\end{lemma}
\begin{proof}
(a) The boundary clause of \cite[Thm.~3.3]{K1} gives $\delta G_j(\square')$ and $\delta C^{(j)}(\square')$ their bounds with the extra factor $e^{-\delta_0\dist(\{u,v\},\partial\square')/\rho_j}$ (the clause used for the lines $\delta G_j(\square',A)$ in the proof of \cref{lem:a02:I} through \cite[(5.3)]{K2}). Writing each factor of the slice $G_{(i)}=a_iG_iQ^*C^{(i)}QG_i$ on $\square'$ as its free-lattice part plus its $\delta$-part and telescoping, every term of $\delta G_{(i)}(\square')$ carries one such factor, and the decay $e^{-\delta_0|\cdot|/\rho_i}$ of the factors (the $\delta$-factor's own included) moves it to both $u$ and $v$ at half the rate ($\min+\text{distance}\ge\max\ge\frac12(\dist(u,\partial\square')+\dist(v,\partial\square'))$), which is $F_i(u)F_i(v)$ with $\delta_1=\delta_0/2$ against the free bounds at the rate $\delta_0/2$; the H\"older clause follows from the sup and first-difference clauses of the factors as \cite[Prop.~3.7]{K1} follows from \cite[Thm.~3.3]{K1}, the boundary factor of the differentiated factor taken at the point of the pair nearer to $\partial\square'$, a pair at distance $>\rho_i$ being bounded by two first differences. On the $f_D$ flow the same follows from the image sum $\delta G_{(i)}(\square')(u,v)=\sum_{g\ne1}\operatorname{sgn}(g)\,G_{(i)}(0)(u,gv)$ of \cref{prop:a02:S}(c), since $|u-gv|\ge\max\{|u-v|,\dist(u,\partial\square')+\dist(v,\partial\square')\}$ for every $g\ne1$ and the images at distance $\le t$ number $O((1+t/p)^d)$. (b) In the ordered bound the line $\ell=(v,v')$ is shrunk by summing one endpoint against its kernel; $\dist(\cdot,\partial\square')$ is $1$-Lipschitz, so $F_i(v)\le e^{\delta_1|v-v'|/\rho_i}F_i(v')$ moves the factor to whichever endpoint survives, at the cost of half of the decay of $\ell$. The surviving vertex $x$ of $T$ is then summed against a leg $l'$ of $T$, a line of index $i'\ge i$ with the kernel profile $|K_{l'}(u-x)|\le C\rho_{i'}^{-c}e^{-\delta_0|u-x|/\rho_{i'}}$ of \cite[(3.63)--(3.65)]{K1} (a cut-line leg, \cref{lem:a02:XL}(b), has the unit-scale profile, $i'=j$), or it is the root. Since $F_i(x)\le\sum_\pi e^{-\delta_1\dist(x,\pi)/\rho_i}$ over the $2d$ face planes of $\square'$, each term depending on one coordinate only, summing that coordinate first gives $\sum_{x_\mu}5^{-j}e^{-\delta_1|x_\mu-\pi_\mu|/\rho_i}\le C\rho_i$ (as $5^{-j}\le\rho_i$), while the sum of $|K_{l'}|$ over the other $d-1$ coordinates is at most $C\rho_{i'}^{-c+d-1}$; so
\[
\sum_{x}5^{-3j}|K_{l'}(u-x)|\,F_i(x)\ \le\ C\,\frac{\rho_i}{\rho_{i'}}\,\rho_{i'}^{\,d-c},
\]
the bound $C\rho_{i'}^{d-c}$ of that step of the induction times $\rho_i/\rho_{i'}$. At the root, where the final position sum is paid by $|\square'|$ (\cite[p.~664]{K1}), $\sum_{x\in\square'}5^{-3j}F_i(x)\le C\rho_i\cdot(\text{area of }\partial\square')\le C(\rho_i/p)\,|\square'|$, the side of $\square'$ being $\ge4p$. In the telescoped product of the ordered bound (\cref{lem:a02:Dsp}(ii)) the factor $\rho_i/\rho_{i'}$ raises by one the exponent of exactly those prefixes that contain $\ell$ and not $l'$, so $D'(T)=D(T)+1=1$ and $D'(\gamma)\in\{D(\gamma),D(\gamma)+1\}$ for $\gamma\supsetneq T$; at the root the credit has no debit.
\end{proof}

\begin{proposition}[the tree at every tuple]\label{prop:a02:TREE}
Let $T$ be the tree \cite[(2.4)]{Ba3} as a prefix component at a fixed tuple, in either branch (with $\square'_K$ in place of $\square'$ in the branch $R>p$), with vertices $v_1$, the reference point $x$, and $v_2$, both summed over $\square'$, its line $\ell$ of index $i$, and the products $W(v_1)$, $W'(v_2)$ of the factors at its vertices (an $\tilde A$ field, the weight $\zeta$ or $1-\zeta$ of step (2) of the proof of \cref{thm:a02:WR}, and one chain-end line of index $\ge i$, a cut line in the branch $R>p$). Then $\sum_{v_1,v_2\in\square'}\ell(v_1,v_2)W(v_1)W'(v_2)$ is a sum of terms of the shapes \textup{(R0)}--\textup{(R2)} of \cref{thm:a02:WR}, with $D'(T)\ge\min\{\alpha,1\}>0$ and $D'(\gamma)\ge D(\gamma)$ for $\gamma\supsetneq T$. In particular the member-wise identity of the tree holds at every tuple, near $\partial\square'$ included, which is what \cref{cor:a02:DEGall}(a) uses for it.
\end{proposition}
\begin{proof}
If $\ell=D_{b_1}D_{b_2}G_{(i)}(0)$, write
\[
\sum_{v_1,v_2\in\square'}\ell\,W\,W'=\sum_{v_1\in\square'}W(v_1)\sum_{v_2\in\square'}\ell(v_1,v_2)\big[W'(v_2)-W'(v_1)\big]+\sum_{v_1\in\square'}W(v_1)W'(v_1)\Big[\sum_{v_2}\ell(v_1,v_2)-\sum_{v_2\notin\square'}\ell(v_1,v_2)\Big].
\]
The full sum over $v_2$ vanishes exactly, a lattice difference summed over the whole lattice (or the torus) being zero, the derivative at the end of a line: \textup{(R0)}. The first term is the transport step of \cite[(3.10)]{Ba3} with $n=0$: the factor $|v_1-v_2|^\alpha$ is absorbed by the decay of $\ell$ at the price $\rho_i^\alpha$, a credit on $\ell$, and the $\alpha$-H\"older quotient of $W'$ costs $\rho_{i'}^{-\alpha}$ on the chain-end line at $v_2$, which is H\"older at its own scale up to $\partial\square'$ (the free part by \cite[(3.65)]{K1}, the $\delta G$ part by \cref{lem:a02:FACE}(a), a cut line by \cref{lem:a02:XL}(b)), and $O(1)$ on the smooth factors $\tilde A$ (\cite[(3.45)]{K1}) and $\zeta$ ($|\nabla\zeta|\le4/p$): \textup{(R2)} with $D'(T)=\alpha$. In the cut term $v_1\in\square'$ and $v_2\notin\square'$, so $|v_1-v_2|\ge\dist(v_1,\partial\square')$ and half of the decay of $\ell$ gives $F_i(v_1)$; the sum over $v_2$ is then bounded by the kernel integral of $\ell$, as at every step of the ordered bound, and \cref{lem:a02:FACE}(b) makes the term an \textup{(R2)} unit with credit $1$. If $\ell=\delta G_{(i)}(\square')$, \cref{lem:a02:FACE}(a) gives $F_i(v_1)$ (the factor $F_i(v_2)$ is dropped) and \cref{lem:a02:FACE}(b) applies; no identity is used. Since $T$ has no partner and its inside family is $\{T\}$, the rewriting is an identity term by term, and the nesting of step (4) of the proof of \cref{thm:a02:WR} is unchanged. Away from $\partial\square'$ and on the torus this is Ba{\l}aban's treatment of \cite[(2.4)]{Ba3} in \cite[Prop.~2.1]{Ba3}; the cut term is the face term that the cut of the position sums produces.
\end{proof}


\paragraph{The cut by clusters.} In the branch $R>p$ King proves \cite[(3.38)]{K1} by treating every propagator longer than $3\sqrt d$ as an external line, so that the graph factorises into localised subgraphs, each reduced with its own cube \cite[p.~666]{K1}. That cut is by vertex positions: at a fixed tuple the partition into subgraphs, hence each subgraph's $\square$, $x_0$, $B$ and $\tilde A$, varies with the positions that the identities of \cref{ssec:a02:ward} sum over, within one member as well as across members. We cut by the centres instead. Units of depth $j$, $d=3$; $c_0\defeq\frac{13}3\sqrt d$ and $R_0\defeq c_0-\frac43\sqrt d=3\sqrt d$, King's threshold. Fix a graph $H$ of $P^{(j),1}$ localised by centres and a tuple $\tau$ with point set $P(\tau)$ (the $y_i$ and the $z_c$).
\begin{enumerate}[label=\textup{(C\arabic*)},nosep,leftmargin=*]
\item \emph{Clusters.} Let $\Gamma_\tau$ be the graph on $P(\tau)$ with an edge between two points at Euclidean distance $\le c_0$; the clusters $K$ are the components of $\Gamma_\tau$, with point sets $P_K$.
\item \emph{Vertices.} A vertex with centres is assigned to the cluster of its centres (well defined on the support of the weights, (a$'$) below); a block or $y$-point carries its own point of $P(\tau)$. An unlocalised vertex is either the attachment vertex of a $\dm$ point vertex whose blob carries an index $\ge j$ (\cite[p.~667]{K1}: no partition of unity is introduced in that graph) or the vertex of a one-vertex graph; the former has exactly two scalar legs \cite[(3.19)]{K1}, and each maximal path through such vertices is read as one \emph{dressed line} between localised ends; one-vertex graphs have no internal line, no divergent subgraph and no cube, and are bounded by the decay of their external lines; an $E_1$ piece localised at one vertex forms one cluster.
\item \emph{Cut.} A line (scalar, photon, external line $a_jG_jQ_j^*$ to a $y$-point, fine $\dm$-blob line or dressed line) is \emph{cut} iff its end-vertices lie in different clusters. A cut line is kept whole, with no scale decomposition, and treated as an external line of both clusters; the \emph{cluster graph} $H_K$ is $H$ restricted to the vertices of $K$, with its uncut lines internal and its cut lines as external fields of index $j$.
\item \emph{Reduction and families.} Each $H_K$ is reduced by \cite[(3.43)--(3.55)]{K1} with its own cube: $\square_K$ the smallest cube of large blocks containing $P_K$, $\square'_K$, $\square''_K$ its $2p$- and $p$-neighbourhoods, $x_0^K$ the centre of $\square_K$, $B_K=A_j(x_0^K)$, $\tilde A_K=A^{(j)}-B_K$ on $\square'_K$; cut lines keep the background $A^{(j)}$ and receive King's unitary factors \cite[(3.46)]{K1} at their ends, as \cite[p.~661]{K1} prescribes for a vertex connected to an external line. Renormalisation families are formed inside each $H_K$ \cite[p.~666]{K1}.
\end{enumerate}
The only change from \cite[p.~666]{K1} is the cut criterion: King decides by vertex positions, (C3) by the clusters of the centres, which are summation variables of \cite[(3.36)]{K1} fixed in $\tau$.

\begin{lemma}[the cluster lemma]\label{lem:a02:CL}
Let $S$ be a prefix component of an ordering of a cluster graph $H_K$ with inside family $F_S$. \textup{(a$'$)} On the support of the localisation weights the centres of one vertex lie in one cluster. \textup{(a)} For $S'\in F_S$ and $H'\defeq(H\setminus S)\cup S'$: $H'$ carries the same tuple under the bijection of \cref{lem:a02:Lsharp}(c), hence the same clusters, cubes $\square_K,\square'_K,\square''_K$, points $x_0^K$, gauges $B_K$, fields $\tilde A_K$ and the same branch; every vertex and line of $S'$ lies in $K$; the legs of $S'$ are the legs of $S$, each with the same cut/uncut status; every line of $H'$ outside $S'$ has its status in $H$. The partition depends on no vertex position. \textup{(b)} On the supports a cut line $u$--$v$ has $|u-v|\ge R_0=3\sqrt d$, and a cut external line from $v$ to the block of a $y$-point has length $>R_0$. \textup{(c)} $\operatorname{diam}P_K\le c_0(|P_K|-1)$, every vertex of $K$ lies within $\frac23\sqrt d$ of $P_K$, and $|\square'_K|\le C(p+c_0N+M)^d\le C'p^d$, $N\defeq|P(\tau)|\le C(\nbar)$. \textup{(d)} The quotient graph (clusters as nodes, cut lines as edges) is connected, and for every spanning tree $T$ of it, $\dist(P(\tau))\le\sum_{\ell\in T}|\ell|+C(N)$, $|\ell|$ the distance between the ends of $\ell$.
\end{lemma}
\begin{proof}
(a$'$) Two centres $c,c'$ of a vertex $v$ have $v\in\supp\chi_{z_c}\cap\supp\chi_{z_{c'}}$, so $|z_c-z_{c'}|\le\frac43\sqrt d<c_0$ and they are adjacent in $\Gamma_\tau$. (a) The clusters are a function of $P(\tau)$, which is common by \cref{lem:a02:Lsharp}(c), and the cluster data are functions of $P_K$; the centres of $S'$ are those of $S$ as points, all in $K$ because $S$ is connected through uncut lines of $H_K$, so by (a$'$) every vertex of $S'$ is in $K$ and every line of $S'$ joins two vertices of $K$; a leg of $S'$ joins a vertex of $S'$ to the same outside end as the corresponding leg of $S$, in some cluster $K_w$, and its status is ``cut iff $K_w\ne K$'' in both; lines outside $S'$ are untouched; no step uses a position. (b) Let $z,z'$ be centres of $u,v$ in different clusters: $|z-z'|>c_0$ and $|u-z|,|v-z'|\le\frac23\sqrt d$. (c) A path in $\Gamma_\tau$ has steps $\le c_0$, and $p\ge1$. (d) Join the points inside each $P_K$ by a tree with steps $\le c_0$ and the clusters along $T$ through the centres of the endpoints of each $\ell$, at the cost $|\ell|+\frac43\sqrt d$ per edge.
\end{proof}
So both of King's phrases, ``treat it as an external line'' and ``the graph then factorizes into a product of localised subgraphs'' \cite[p.~666]{K1}, become exact, with the same factorisation for every member and every position of every family.

\begin{lemma}[regularised cut lines]\label{lem:a02:XL}
Fix a radial $\bar\psi\in C^\infty(\R^d)$, $0\le\bar\psi\le1$, $\bar\psi(w)=0$ for $|w|\le\frac12$, $\bar\psi(w)=1$ for $|w|\ge1$, and for a cut line $\ell$ between $u\in K$ and $v\in K'$ put $\ell^\flat(u,v)\defeq\bar\psi(u-v)\ell(u,v)$ (for an external line to a $y$-point, $\bar\psi(u-y)$), and $\ell^{\flat,\mathrm{ph}}\defeq U(B_K(\Gamma_{x_0^K,u}))^{-1}\ell^\flat U(B_{K'}(\Gamma_{x_0^{K'},v}))$ for the phase-dressed form of (C4). \textup{(a)} Replacing every cut line by $\ell^\flat$ leaves $E^{(j)}(H;\tau)$ unchanged, term by term: for every member of every family, every ordering and every position configuration, each evaluation pair $(u',v')$ at which a cut line enters, the lattice differences of the vertex functions at its ends included, has $|u'-v'|\ge R_0-4\cdot5^{-j}>1$, where $\bar\psi=1$. \textup{(b)} $\ell^{\flat,\mathrm{ph}}$ is $C^{1,\alpha}\otimes C^{1,\alpha}$ at unit scale in both arguments with decay $e^{-\delta_1|u-v|}$, at all positions and with no distance condition, the constant depending on $\nbar$, $d$, $L$, $\alpha$ only, times $(1+k-j)^5$ for a dressed line (the size of the coarse $\dm$ blob); this is \cite[(3.1)]{Ba3} in the form of a cut line with the endpoint localisations.
\end{lemma}
\begin{proof}
(a) The ends of a cut line are localised in the original sum: (C2) assigns to a cluster localised vertices, whose weights are present, block and $y$-points, which are points of $P(\tau)$, the vertices of one-centre $E_1$ pieces, none of whose lines is cut, and unlocalised coarse $\dm$ vertices, which lie strictly inside dressed lines and are never an end; so \cref{lem:a02:CL}(b) applies at every evaluation pair up to the lattice differences, each of which moves an end by at most $5^{-j}$. (b) For $|u-v|\ge\frac12$ the slice bounds \cite[Thm.~3.3, Props.~3.7--3.9]{K1} (\cref{prop:a02:FL}, \cref{prop:a02:S}) of the whole line, its derivatives and its mixed second differences are those of a unit-scale kernel with the stated decay, and the phases \cite[(3.46)]{K1} and $\bar\psi$ are smooth at unit scale; for $|u-v|<\frac12$ the kernel vanishes.
\end{proof}

\begin{theorem}[the Ward rows at every tuple]\label{thm:a02:WRall}
Let the localisation be by centres and the cut by clusters, with every cut line replaced by $\ell^\flat$, and let $S$ be a Ward-row prefix component of an ordering of a cluster graph $H_K$ at a fixed tuple $\tau$ (in particular in the branch $R>p$), with $A'$-line legs and scalar legs that may be cut lines. Then the sum over $F_S$ and $A_S$ is a sum of terms of the shapes of \cref{thm:a02:WR} with $\square\to\square_K$: \textup{(R0)} zero; \textup{(R1)} $\sum_x5^{-3j}c_a(x)\prod\partial^a\tilde\Phi(x)$, $|a|\le-D(S)$, $|c_a|\le C\rho_{j_0}^{D(S)+|a|}$, times $W^\circ$ and the frozen $\tilde A_K(x)$-legs, a cut-line leg counting as index $j$; \textup{(R2)} a credit $\rho^\theta$ on a line of $S$ with a debit of order $\le\theta$ on one leg (a cut-line leg costs $1$, by \cref{lem:a02:XL}(b)) or on a unit-scale weight, $D'(S)>0$, $D'(\gamma)\ge D(\gamma)$ for $\gamma\subsetneq S$. $C$ is uniform in $x$, $j_0$, $k$, $\eta$ and $\tau$ and depends on $\nbar$, $d$, $L$, $\alpha$ only, up to the dressed-line factor $(1+k-j)^5$ of \cref{lem:a02:XL}(b), absorbed in $(1+k)^{p_0}$ of \cref{cor:a02:DEGall}(b); the decay budget is $\nbar+4$ draws per line, and every evaluation point of a cut-line end is joined to a localised anchor of its own cluster by at most $C(\nbar)$ lines of rewritten units.
\end{theorem}
\begin{proof}
The proof of \cref{thm:a02:WR} is run inside $H_K$ with these changes. By \cref{lem:a02:CL}(a) every member of every family met in the nesting is reduced with the same cluster data, and cut lines never lie inside a unit: no ordered subgraph contains a cut line (it carries index $j$ and is external to both clusters), so a loop through a cut line is never a divergent unit. The frozen-weight identity of \cref{lem:a02:FW} holds per cluster, the legs that are cut lines entering through $\ell^{\flat,\mathrm{ph}}$, which by \cref{lem:a02:XL}(b) is a $C^{1,\alpha}$ leg at unit scale, so its debit in the Taylor and transport steps costs $1$; the identities of \cref{lem:a02:GC,lem:a02:REF} act only on $F^C_S$ after \cref{lem:a02:RS}, inside one cluster, and the cut lines are legs that enter only through the transport and Taylor steps after the identity, so the cut adds no identity and truncates no member line. The cut-line factor of an output term is bounded by $Ce^{-\delta_2d(K,K')}$ times transfer factors on unit lines: each evaluation point of a cut-line end is joined to a localised anchor by at most $C(\nbar)$ lines of its own cluster's rewritten units, the transfer takes one more decay draw per line, and the bound factorises over clusters. Nesting is as in \cref{thm:a02:WR}(4), with $\square_K$ in place of $\square$.
\end{proof}

\begin{corollary}[the degree count on every row; the bound in the branch $R>p$]\label{cor:a02:DEGall}
\textup{(a)} With the localisation by centres and the cut by clusters, every source-free class with $D\le0$ obeys \DEG\ at every tuple: the non-Ward rows by their member-wise identities (charge parity, the derivative at a line end, reflection, the bounded $\eta$-suppressed tadpole) or by the counterterm matching between partners of identical vertex structure, whose centre multisets coincide by \cref{lem:a02:Lsharp}(b) (the only non-Ward $D=-1$ class with a Taylor step is $\Sigma_1$, whose two centres sit at its two vertices, the attachment points of its scalar legs, so its non-reference centre is folded into the finer leg, \cite[(3.85)]{K1}, exactly as under King's rule); the Ward rows by \cref{thm:a02:WR,thm:a02:WRall}. \textup{(b)} For every graph whose divergent units are such classes, $|E^{(\theta)}(H_{\mathrm{ren}};\tau)|\le C\eta^{\rho_0}p^{n_1}(1+k)^{p_0}\exp[-\delta\dist(P(\tau))]$, $\theta\in\{k,k+n\}$, at every tuple, which is \cite[(3.38)]{K1} in the branch $R>p$; with \cite[(3.39)]{K1} from \cref{thm:a02:WR} in the branch $R\le p$ this gives King's per-graph estimate \cite[(3.42)]{K1}. \textup{(c)} At tuples with at least two clusters every connected component of every cluster graph carries a cut-line leg, so the vacuum members and the component-wise cancellation of \cite[p.~667]{K1} do not arise there, and $E(H;A^{(j)};\tau)$ and $E(H;0;\tau)$ are bounded separately; a closed scalar loop of internal lines of one cluster graph none of whose vertices carries another charged field vanishes when the total $q$-power on it is odd (no phase sits on the loop), which is the per-cluster form of charge parity.
\end{corollary}
\begin{proof}
(a) is the list of mechanisms of \cref{ssec:a02:ward} with the two theorems. (b) Per cluster, \cref{thm:a02:WRall} and (a) give the rewritten form with every stage of positive degree; the multiscale bound \cref{thm:a02:S3}(a) with the root held, the cut-line factors of \cref{lem:a02:XL} and \cref{lem:a02:CL}(d) (the tree distance of the tuple is controlled by the cut lines of a spanning tree) give the exponential in $\dist(P(\tau))$, and the volume factors $|\square'_K|\le C'p^d$ of \cref{lem:a02:CL}(c) give $p^{n_1}$. (c) A cluster graph is connected to the rest of $H$ through cut lines only, and $H$ is connected; on a closed loop with no charged external field the internal phases cancel and $\Tr q^{2r+1}=0$.
\end{proof}

For \cref{thm:a02:T} this supplies the $f=1$ input of (ii) on the Ward rows, $|c_a^1|\le C\rho_{j_0}^{D(S)+|a|}$, at every tuple of either branch; the transfer to the $f_D$ flow is (i)--(iii) there, the exact-zero and Ward mechanisms being identities for every charge-diagonal line kernel when partners carry the same lines, and \cref{lem:a02:Lsharp}(b) holding for the $f_D$ members verbatim since the weights do not enter the centre count.
